\chapter{Cavernous Sinus Thrombosis}

\section*{Definition and Overview}
Cavernous sinus thrombosis is a rare but life-threatening condition in which a blood clot (thrombus) forms in the cavernous sinus, a pair of large veins located at the base of the brain, just behind each eye. The cavernous sinus drains blood from the face, eyes, and brain, and it carries several important nerves that control eye movement, sensation over the face, and the blink reflex.

The condition usually develops as a complication of an infection spreading from the face, nose, sinuses, teeth, or ears, and it causes severe headache, eye pain, swelling, and paralysis of the eye muscles. It is a medical emergency that requires urgent hospital treatment with antibiotics and blood-thinning medication. Prompt diagnosis and treatment have greatly improved outcomes, but the condition still carries a real risk of death or permanent disability.

\section*{Epidemiology}
Cavernous sinus thrombosis is rare, but it is one of the most feared complications of infections of the face and paranasal sinuses. It most often follows staphylococcal infection of the face, particularly boils and infections around the nose and upper lip, or spread of infection from the sinus cavities, middle ear, or teeth. It can occur at any age but is most common in younger adults, and it develops more often when facial infections are squeezed, picked, or left untreated.

Because it is so uncommon, its exact frequency is difficult to measure, and figures rely on case reports and case series. Mortality has fallen substantially since the introduction of antibiotics, from levels above 80\% in the pre-antibiotic era to single figures in modern series, but permanent complications remain common among survivors.

\section*{Aetiology and Pathophysiology}
The cavernous sinus is particularly vulnerable to infection because of its anatomy. Veins from the central face, nose, and sinuses drain backward into it, and unlike most veins, they contain no valves, so blood (and bacteria) can flow in either direction. The steps are:

\begin{itemize}
    \item \textbf{Source infection:} usually a boil or cellulitis of the face, especially around the nose and upper lip, or sinusitis, dental infection, or middle ear infection
    \item \textbf{Spread:} bacteria travel along the facial and ophthalmic veins into the cavernous sinus
    \item \textbf{Inflammation and clotting:} the infection inflames the sinus lining and triggers clot formation
    \item \textbf{Local pressure:} the swollen, clotted sinus presses on the nerves and structures around it, and raised pressure impairs drainage of the eye and face
    \item \textbf{Systemic spread:} in severe cases, bacteria enter the bloodstream, and the infection can spread within the skull to the meninges or brain
\end{itemize}

Staphylococcus aureus is the most commonly identified organism, though streptococci, pneumococci, and anaerobic bacteria also occur.

\section*{Clinical Features}
\begin{itemize}
    \item Severe headache, usually on one side at first and often behind or around the eye
    \item High fever, chills, and a feeling of being severely unwell
    \item Swelling, redness, and protrusion of the eye (proptosis)
    \item Pain on moving the eye, double vision, and drooping of the eyelid
    \item Difficulty moving the eye in various directions (ophthalmoplegia)
    \item Facial pain or loss of sensation over the forehead and cheek
    \item Redness and swelling over the bridge of the nose and eyelids
    \item Symptoms often spread to involve both eyes within a day or two, because the sinuses connect
    \item Nausea, vomiting, and in severe cases confusion, seizures, or loss of consciousness
\end{itemize}

The illness usually follows a facial or sinus infection by days, and the onset of eye symptoms is often sudden.

\section*{Differential Diagnosis}
\begin{itemize}
    \item \textbf{Orbital cellulitis:} infection of the tissues of the eye socket, which can produce similar swelling and pain but without the typical nerve involvement
    \item \textbf{Severe sinusitis:} sinus infection alone may cause facial pain, but not the eye findings of cavernous sinus thrombosis
    \item \textbf{Migraine or cluster headache} with eye symptoms
    \item \textbf{Orbital or facial tumours} pressing on the eye and its nerves
    \item \textbf{Other causes of sudden double vision:} such as diabetes-related nerve palsy or a stroke affecting eye movement
    \item \textbf{Carotid-cavernous fistula:} an abnormal connection between an artery and the cavernous sinus, often following head injury
\end{itemize}

The combination of a facial or sinus infection with headache, fever, and eye signs points strongly to cavernous sinus thrombosis, but imaging is required to confirm the diagnosis.

\section*{When to Seek Medical Care}
Seek urgent medical attention for any severe headache, fever, or swelling of the face or eye, especially if there is pain on moving the eye, double vision, or drooping of the eyelid, and particularly within days of a boil, facial infection, or bout of sinusitis.

\textbf{Call 000 (or local emergency services) for:} any of the features above with confusion, seizures, or reduced consciousness, or a rapidly worsening illness. Cavernous sinus thrombosis is a medical emergency and must be assessed in hospital without delay.

\section*{Diagnosis and Investigations}
\begin{itemize}
    \item \textbf{History and examination:} identification of the source infection, the onset and spread of symptoms, and a full neurological and eye examination, including eye movements, pupils, and sensation over the face
    \item \textbf{Imaging:} MRI with venography is the most sensitive investigation, and CT with venography is a fast alternative in the emergency setting
    \item \textbf{Lumbar puncture:} examination of the cerebrospinal fluid may be needed to exclude meningitis, though it is less useful in confirming the clot itself
    \item \textbf{Blood tests:} full blood count, inflammatory markers, and blood cultures to identify the organism
    \item \textbf{Ear, nose, and throat assessment:} examination of the sinuses, ears, and teeth to find and treat the source of infection
\end{itemize}

\section*{Management}
\begin{itemize}
    \item \textbf{Hospital admission:} typically to a high-dependency or intensive care setting for close monitoring
    \item \textbf{Intravenous antibiotics:} high-dose antibiotics covering the likely organisms, usually including anti-staphylococcal cover, given for several weeks
    \item \textbf{Anticoagulation:} blood-thinning treatment with heparin to limit clot extension, guided by the clinical team and the risks of bleeding
    \item \textbf{Treatment of the source:} drainage of the infected sinus, drainage of an abscess, or removal of an infected tooth
    \item \textbf{Supportive care:} fluids, pain relief, control of fever, and monitoring of eye pressure and vision
    \item \textbf{Surgery:} occasionally needed to drain an abscess or to relieve pressure on the optic nerve
    \item \textbf{Follow-up:} prolonged ophthalmology and neurology follow-up, because complications can appear later
\end{itemize}

\section*{Complications}
\begin{itemize}
    \item Meningitis or brain abscess from spread of infection within the skull
    \item Septicemia and septic shock
    \item Damage to the nerves controlling eye movement, causing permanent double vision
    \item Reduced vision or blindness from pressure on the optic nerve or impaired blood supply to the eye
    \item Pituitary gland damage, causing hormonal deficiency
    \item Stroke-like damage from impaired venous drainage of the brain
    \item Death in a small proportion of cases, despite treatment
\end{itemize}

\section*{Prognosis and Outcomes}
With prompt antibiotic treatment and anticoagulation, the majority of people now survive cavernous sinus thrombosis, but recovery is often slow and incomplete. Headaches and double vision may persist for many months, and a significant proportion of survivors are left with permanent neurological deficits, including double vision, facial weakness, or visual loss. Early recognition and treatment of the source infection are the strongest influences on outcome, and untreated or delayed cases carry a high risk of death.

\section*{Prevention and Self-Care}
\begin{itemize}
    \item Do not squeeze, pick, or lance boils or pimples, particularly around the nose, lips, and central face
    \item Treat facial infections and cellulitis promptly with appropriate care
    \item Manage sinusitis, dental infections, and ear infections early, especially if they are painful or persistent
    \item Keep up to date with dental care to prevent dental abscesses
    \item Seek medical advice promptly for any spreading facial infection, particularly with fever
    \item After treatment, attend all follow-up appointments for eye and nerve assessment
\end{itemize}

\section*{Sources and Further Reading}
The following organisations provide reliable information on cavernous sinus thrombosis and related emergencies.

\begin{itemize}
    \item National Health Service. \textit{NHS conditions guide on head, brain, and eye emergencies.} https://www.nhs.uk/conditions/
    \item Mayo Clinic. \textit{Information on sinus infections and their complications.} https://www.mayoclinic.org/
    \item MSD Manual. \textit{Cavernous sinus thrombosis -- professional overview.} https://www.msdmanuals.com/
    \item MedlinePlus. \textit{Consumer health information on brain and nervous system conditions.} https://medlineplus.gov/
    \item Centers for Disease Control and Prevention. \textit{Blood clotting and stroke information.} https://www.cdc.gov/
    \item World Health Organization. \textit{Global guidance on neurological and infectious diseases.} https://www.who.int/
\end{itemize}
